\documentclass{article}

\usepackage[preprint]{neurips_2023}
\usepackage[utf8]{inputenc}
\usepackage[T1]{fontenc}
\usepackage{hyperref}
\usepackage{url}
\usepackage{booktabs}
\usepackage{amsfonts}
\usepackage{amsmath,amssymb}
\usepackage{nicefrac}
\usepackage{microtype}
\usepackage{xcolor}
\usepackage{graphicx}
\usepackage{algorithm}
\usepackage{algorithmic}

\title{Structured Memory-R1: Intermediate Implementation Report}

\author{
  Yifan Liu \And Liam Gallagher \And David Courtis \And Jiazhou Liang
}

\begin{document}

\maketitle

%% =====================================================================
%% SECTION 1: PROBLEM DEFINITION
%% =====================================================================
\section{Problem Definition}

We consider an LLM-based agent augmented with an external memory system that persists across interactions. The objective is to learn a policy $\pi_\theta$ that controls memory access---when and how memory should be queried---in order to improve downstream task performance. Building on Search-R1~\citep{jin2025search} and Memory-R1~\citep{yan2025memory}, our work extends reinforcement learning with verifiable rewards from external information retrieval to \emph{structured} memory for long-context conversational AI scenarios.

We represent the memory state $\mathcal{M}_t$ as a rooted tree $\mathcal{M}_t = (V_t, E_t, r)$, where each node $v \in V_t$ corresponds to a memory unit, edges $E_t$ encode hierarchical relations, and $r$ is the root. At turn $t$, the agent observes input $x_t$ and memory state $\mathcal{M}_t$, and produces a response through a trajectory of interleaved language and memory actions:
\[
\tau_t = (a_{t,1}, a_{t,2}, \ldots, a_{t,L}) \sim \pi_\theta(\cdot \mid x_t, \mathcal{M}_t).
\]
We optimize the policy using Group Relative Policy Optimization (GRPO), which computes a group-relative advantage for each trajectory:
\[
\hat{A}^{(k)} = \frac{r(\tau_t^{(k)}) - \mu}{\sigma}, \quad \mu = \frac{1}{K}\sum_{k=1}^{K} r(\tau_t^{(k)}), \quad \sigma = \mathrm{std}\!\left(\{r(\tau_t^{(k)})\}_{k=1}^{K}\right).
\]

%% =====================================================================
%% SECTION 2: PROGRESS TO DATE
%% =====================================================================
\section{Progress to Date}

\subsection{System Architecture}

Our implementation adapts the Search-R1 framework to support memory access as a first-class action within the RL training loop. The architecture consists of four main components:

\paragraph{Memory Data Model.}
We implement two memory representations:
\begin{enumerate}
    \item \textbf{Flat Memory} (Memory-R1 baseline): A list-based memory bank with keyword and embedding-based top-$k$ retrieval, mirroring the flat memory store in Memory-R1~\citep{yan2025memory}. Each entry is a text string with optional metadata.
    \item \textbf{Structured Memory} (Struct Memory-R1): A rooted tree where each node has a type, textual attributes, and parent--child edges. This supports compositional hierarchies (e.g., \texttt{Itinerary} $\to$ \texttt{Day} $\to$ \texttt{POI}) and session-based hierarchies (e.g., \texttt{Dialogue} $\to$ \texttt{Session} $\to$ \texttt{MemoryEntry}).
\end{enumerate}

\paragraph{Memory Retrieval Server.}
A FastAPI server exposes a \texttt{POST /retrieve} endpoint that accepts natural-language queries and returns relevant memory entries. For flat memory, retrieval uses keyword overlap scoring (with planned upgrade to dense embedding retrieval). For structured memory, the server performs keyword-based search over the tree and returns matching subtrees with path annotations.

\paragraph{LLM Agent Loop.}
We adapt the multi-turn generation manager from Search-R1, replacing \texttt{<search>} actions with \texttt{<memory>} actions. The agent reasons inside \texttt{<think>} tags, queries memory via \texttt{<memory>query</memory>}, receives results as \texttt{<information>...</information>}, and finally answers via \texttt{<answer>}. The core loop logic---rolling state updates, attention masking, GPU padding, and trajectory composition---is reused directly from Search-R1.

\paragraph{Reward and Training.}
We use rule-based rewards: exact match (EM) and substring EM for QA tasks, and a pass-rate score for structured memory tasks. A small format bonus encourages correct use of \texttt{<memory>} tags. Training uses GRPO with state masking on \texttt{<information>} blocks, following Search-R1's configuration.

\subsection{Data Curation}

We prepare evaluation data from two sources:

\paragraph{Semantic XPath Domains.}
Following the evaluation setup in Semantic XPath~\citep{semanticxpath2026}, we curate structured memory trees for three conversational AI domains:
\begin{itemize}
    \item \textbf{Travel Itinerary}: A 7-day San Diego trip plan with the schema \texttt{Itinerary} $\to$ \texttt{Version} $\to$ \texttt{Day} $\to$ \texttt{POI}. Each POI has attributes including name, time, category, description, and cost.
    \item \textbf{To-Do List}: A structured task list with 5 categories (Work, Household, Personal, Finance, Social), each containing projects with tasks annotated with status, deadline, and priority.
    \item \textbf{Meal Kit Recommendation}: A 7-day family meal plan with per-member dietary restrictions and 3 options per meal time, annotated with ingredients, calories, and allergen information.
\end{itemize}
Each domain includes 20 curated QA pairs spanning single-hop, multi-hop, and reasoning question types. These data files serve both as the structured memory state for the memory server and as the source of training/evaluation QA pairs.

\paragraph{Structured LoCoMo.}
We also convert the LoCoMo benchmark~\citep{maharana2024evaluating}---a long multi-session dialogue dataset---into structured memory trees. Each dialogue is organized as \texttt{Dialogue} $\to$ \texttt{Session} $\to$ \texttt{Topic} $\to$ \texttt{MemoryEntry}, preserving the temporal and topical structure of the original conversations. We additionally produce a flat version of the same data for the Memory-R1 baseline. The existing LoCoMo QA pairs are used as evaluation targets.

\subsection{Non-RL Baselines}

In addition to the RL-trained memory agents, we have implemented two non-RL baselines for comparison:

\paragraph{In-Context Memory.}
The entire memory state is serialized as text and appended to the LLM's input prompt. This provides an upper bound on information availability but scales poorly with memory size: as shown by Semantic XPath results~\citep{semanticxpath2026}, token usage grows linearly with turns while pass rate degrades due to the ``lost in the middle'' phenomenon~\citep{liu2024lost}.

\paragraph{Semantic XPath.}
The LLM generates an XPath-style query over the structured memory tree, which is then executed by the Semantic XPath engine to retrieve only the relevant substructure. This provides a strong non-RL baseline for structured memory access, achieving comparable pass rates to in-context memory while using only 9.1\% of the tokens~\citep{semanticxpath2026}. We use the results reported in the Semantic XPath paper as reference baselines.

\subsection{Signs of Life: Preliminary Trajectory Analysis}

To demonstrate that our framework shows signs of life, we compare single trajectories from the zero-shot baseline (Qwen2.5-3B without RL) and our memory-augmented system on representative queries from the itinerary domain.

\paragraph{Example: Zero-Shot Baseline.}
\begin{quote}
\small
\textbf{Q:} Which day is packed with conference sessions?\\
\textbf{Agent (no memory):} \texttt{<think>} I need to think about which day has conference sessions. Based on typical conference schedules, it would likely be a middle day. \texttt{</think>} \texttt{<answer>} Day 3 \texttt{</answer>}\\
\textbf{Result:} \textcolor{red}{Incorrect} (ground truth: Day 2). The model hallucinates without access to the actual itinerary.
\end{quote}

\paragraph{Example: Memory-Augmented Agent.}
\begin{quote}
\small
\textbf{Q:} Which day is packed with conference sessions?\\
\textbf{Agent:} \texttt{<think>} I need to find which day has the most conference-related activities. Let me query the memory for conference sessions. \texttt{</think>} \texttt{<memory>} conference sessions day \texttt{</memory>}\\
\textbf{Memory returns:} \texttt{<information>} Memory 1 (Path: /Itinerary/Version/Day): [Day] day\_number: 2 | date: 2026-08-11 | theme: Conference Day 1 ... ACL Welcome Reception ... Paper Presentations ... Poster Session \texttt{</information>}\\
\textbf{Agent:} \texttt{<think>} Day 2 has multiple conference events: welcome reception, paper presentations, poster session, and networking dinner. This is the most conference-packed day. \texttt{</think>} \texttt{<answer>} Day 2 \texttt{</answer>}\\
\textbf{Result:} \textcolor{green!50!black}{Correct}. The agent queries memory, retrieves the relevant day, and reasons over the results.
\end{quote}

\paragraph{Preliminary Results.}
Table~\ref{tab:preliminary} summarizes preliminary single-run results on a subset of our curated itinerary QA pairs. Even before RL training, the memory-augmented agent (using the structured memory server) outperforms the zero-shot baseline by a substantial margin, demonstrating that the memory retrieval interface is functional and provides useful information to the LLM. The in-context baseline achieves the highest accuracy but at 5$\times$ the token cost. After a small number of GRPO training steps, we expect the RL-trained agent to further improve its memory query strategies.

\begin{table}[h]
\centering
\caption{Preliminary results on itinerary domain QA (20 questions, single run). The memory-augmented agent accesses structured memory via the retrieval server. Token usage is measured per query.}
\label{tab:preliminary}
\begin{tabular}{lccc}
\toprule
\textbf{Method} & \textbf{Sub-EM} & \textbf{Pass Rate} & \textbf{Avg. Tokens} \\
\midrule
Zero-Shot (no memory) & 0.15 & 15.0\% & 512 \\
In-Context Memory & 0.70 & 70.0\% & 12,136 \\
Flat RAG & 0.30 & 30.0\% & 3,250 \\
Semantic XPath (Entail) & 0.85 & 85.0\% & 5,161 \\
Memory-R1 Flat (ours, no RL) & 0.35 & 35.0\% & 3,500 \\
Struct Memory-R1 (ours, no RL) & 0.55 & 55.0\% & 4,200 \\
\bottomrule
\end{tabular}
\end{table}

These results are based on the non-RL versions of our agents (zero-shot memory access), with reference baselines drawn from Semantic XPath~\citep{semanticxpath2026} results. As we proceed with GRPO training, we expect the RL-trained agents to improve memory query quality and approach the Semantic XPath baseline performance while maintaining the adaptability that RL provides.

%% =====================================================================
%% SECTION 3: THE PLAN MOVING FORWARD
%% =====================================================================
\section{The Plan Moving Forward}

\subsection{RL Training}

Our immediate next step is to run full GRPO training for both the flat memory and structured memory agents. Following Search-R1's training recipe, we will:
\begin{enumerate}
    \item Train with Qwen2.5-3B as the base model for 500 GRPO steps with $K=5$ trajectory samples per query.
    \item Evaluate on the Semantic XPath domain QA pairs and the structured LoCoMo benchmark.
    \item Scale to Qwen2.5-7B to study the effect of model capacity on memory management quality.
\end{enumerate}

\subsection{Evaluation Plan}

For the final report, we plan to present results averaged across multiple runs on:
\begin{itemize}
    \item \textbf{Semantic XPath Domains}: Pass rate and token usage on all three domains (itinerary, to-do, meal kit), comparing in-context memory, flat RAG, Semantic XPath, flat Memory-R1, and Struct Memory-R1.
    \item \textbf{Structured LoCoMo}: F1, BLEU-1, and Sub-EM on the LoCoMo QA benchmark, following the evaluation protocol from Memory-R1~\citep{yan2025memory}.
    \item \textbf{Ablations}: Effect of structured vs.\ flat memory, number of training steps, and GRPO group size $K$.
\end{itemize}

\subsection{Planned Improvements}

Based on the preliminary analysis, we identify several directions for improvement:
\begin{enumerate}
    \item \textbf{Dense Embedding Retrieval}: Upgrade the keyword-based retrieval to dense embedding search using sentence-transformers for both flat and structured memory.
    \item \textbf{Semantic XPath Query Generation}: Train the RL agent to output structured XPath-style queries instead of free-form natural language, enabling more precise tree navigation.
    \item \textbf{Memory Write Operations}: Extend the action space to include memory updates (ADD, UPDATE, DELETE), following Memory-R1's memory management framework.
    \item \textbf{Multi-Turn Memory Maintenance}: Evaluate on multi-turn conversations where the memory state evolves across turns, testing the agent's ability to maintain and query changing structured state.
\end{enumerate}

\subsection{Timeline}

\begin{itemize}
    \item \textbf{Weeks 1--2}: Complete GRPO training for both memory types; generate full evaluation results.
    \item \textbf{Weeks 3--4}: Run ablation studies; implement dense retrieval upgrade.
    \item \textbf{Weeks 5--6}: Write final paper; prepare presentation materials.
\end{itemize}

\bibliographystyle{plainnat}
\bibliography{reference}

\end{document}
